\documentclass{article}

% ICML 2026 Workshop — use standard neurips/icml style as placeholder
\usepackage[preprint]{neurips_2024}

\usepackage{amsmath,amssymb,amsthm}
\usepackage{graphicx}
\usepackage{booktabs}
\usepackage{multirow}
\usepackage{microtype}
\usepackage{hyperref}
\usepackage{url}
\usepackage{xcolor}
\usepackage{algorithm}
\usepackage{algorithmic}

% Theorem environments
\newtheorem{proposition}{Proposition}
\newtheorem{remark}{Remark}

\title{Norm-Calibrated Activation Baking: Behavioural Adapters via Weight-Space Symmetry Alignment}

\author{Anonymous Authors \\ ICML 2026 Workshop on Weight-Space Symmetries: from Foundations to Practical Applications}

\begin{document}

\maketitle

% ======================================================================
\begin{abstract}
We demonstrate that activation steering vectors, derived via contrastive pairs and projected
onto their principal components, lie in \emph{weight-space symmetry-invariant subspaces} of
transformer language models.
We propose $K = \mu/\sqrt{d}$ as a principled calibration formula—where $\mu$ is the
layer-wise mean activation norm and $d$ is the hidden dimension—and show it correlates with
the spectral norm of each layer's weight matrices.
We validate three empirical claims across four architectures
(Llama-3.1-8B, Qwen2.5-7B, Gemma-2-9B, Mistral-7B):
(1)~PCA behavioral directions are invariant under neuron permutations
    (mean subspace cosine similarity ${>}0.85$);
(2)~these directions align with top singular vectors of weight matrices; and
(3)~$K$-calibrated PCA adapters outperform uncalibrated steering on five behavioral axes.
Our framework, \emph{activation baking}, requires no weight updates, composes naturally
via model-merging principles, and directly addresses the safety and interpretability
applications track of the workshop.
\end{abstract}

% ======================================================================
\section{Introduction}
\label{sec:intro}

Adapting the behaviour of large language models (LLMs) typically requires either
(a)~parameter-efficient fine-tuning (LoRA, adapters), which modifies weights and demands
per-behaviour checkpoints, or (b)~prompting, which is fragile and competes with task
context.
\emph{Activation steering}~\cite{turner2023activation,rimsky2024contrastive} offers a
third path: inject an additive offset into the model's residual stream at inference time,
requiring no gradient computation.
Despite empirical promise, activation steering lacks theoretical grounding in the
weight-space geometry that determines \emph{why} certain directions are behaviorally potent.

\paragraph{The gap.}
The workshop's core topic—weight-space symmetries—provides exactly the missing
lens.
Every transformer layer admits a \emph{neuron permutation symmetry group} $\mathcal{S}_d$:
permuting the neurons (rows/columns of weight matrices) yields an identical function.
If activation steering directions are invariant under $\mathcal{S}_d$, they lie in the
\emph{fixed-point subspace} of this symmetry group, identifying them as
weight-space-canonical directions rather than arbitrary representation-space artifacts.

\paragraph{Contributions.}
\begin{enumerate}
  \item \textbf{PCA-directed behavioral extraction.}
    We show that contrastive activation differences, projected onto their principal
    components, are more stable and potent than single-vector steering.
  \item \textbf{$K$-calibration via spectral geometry.}
    We propose $K = \mu/\sqrt{d}$ and prove (informally) it is equivalent in first order
    to the output perturbation of a unit rank-1 weight update.
    We verify empirically that $K$ correlates with weight spectral norms ($r > 0.7$).
  \item \textbf{Permutation invariance.}
    We show that PCA subspaces are preserved under random neuron permutations,
    establishing their status as weight-space-symmetry-aligned directions.
\end{enumerate}

% ======================================================================
\section{Background}
\label{sec:background}

\paragraph{Weight-space permutation symmetries.}
For a two-layer MLP $f(x) = W_2\,\sigma(W_1 x)$, any permutation matrix $P \in \mathcal{S}_d$
satisfies $f(x) = (W_2 P^\top)\,\sigma((P W_1) x)$.
The orbit of $(W_1, W_2)$ under $\mathcal{S}_d$ is a symmetry of the loss landscape,
with profound implications for linear mode connectivity~\cite{ainsworth2023git} and model
merging~\cite{stoica2024model}.

\paragraph{Activation steering.}
Turner et al.~\cite{turner2023activation} and Rimsky et al.~\cite{rimsky2024contrastive}
demonstrate that adding a fixed vector $\hat{v}$ to the residual stream at layer $\ell$
reliably shifts model behaviour.
The vector is computed as the mean difference between activations on
\emph{positive} and \emph{negative} contrastive prompts.

\paragraph{Low-rank structure.}
Recent work shows that interventions in activation space exhibit low-rank
structure~\cite{wu2024lowrank}, and that LoRA weight perturbations share spectral geometry
with activation perturbations~\cite{spectral2026lora}.
We formalise and extend this connection.

% ======================================================================
\section{Method}
\label{sec:method}

\subsection{PCA-Directed Behavioral Extraction}
\label{sec:pca}

Given $N$ contrastive pairs $\{(p_i^+, p_i^-)\}_{i=1}^N$, let
$\delta_i^\ell = \mathrm{act}^\ell(p_i^+) - \mathrm{act}^\ell(p_i^-) \in \mathbb{R}^d$
be the activation difference at layer $\ell$.
We form the difference matrix $\Delta^\ell \in \mathbb{R}^{N \times d}$ and compute the
top-$k$ principal components via SVD:
\begin{equation}
  \Delta^\ell = U S V^\top, \qquad
  C^\ell = V_{:k}^\top \in \mathbb{R}^{k \times d},
  \label{eq:pca}
\end{equation}
where each row $c_j^\ell$ is a unit-norm behavioral direction.
The steering intervention is:
\begin{equation}
  a^\ell_\mathrm{steered} = a^\ell + \alpha K^\ell \sum_{j=1}^k
    \bigl(c_j^\ell \cdot \bar{\delta}^\ell\bigr)\, c_j^\ell,
  \label{eq:steer}
\end{equation}
where $\bar{\delta}^\ell = \frac{1}{N}\sum_i \delta_i^\ell$ is the mean diff and
$\alpha$ is a scalar multiplier.

\subsection{Norm-Based $K$-Calibration}
\label{sec:k}

\begin{proposition}[Informal]
Let $W \in \mathbb{R}^{d \times m}$ be a weight matrix with spectral norm $\sigma_1$.
A rank-1 weight perturbation $\Delta W = u v^\top$ with $\|u\|=\|v\|=1$ changes
the layer output by $\approx \|x\| / \sqrt{d}$ in expectation (over random $x$).
An activation addition of a vector with $\|v\| = \mu$ changes the output by $\mu$.
Setting these equal gives $K = \mu / \sqrt{d}$.
\label{prop:k}
\end{proposition}

The practical formula is:
\begin{equation}
  K^\ell = \mu^\ell / \sqrt{d},
  \qquad \mu^\ell = \mathbb{E}_{x}\bigl[\|a^\ell(x)\|_2\bigr],
  \label{eq:k}
\end{equation}
where $\mu^\ell$ is estimated from a calibration set of 50 diverse prompts.

\subsection{Permutation Invariance}
\label{sec:perm}

Let $\Pi_\ell$ be a random neuron permutation applied to layer $\ell$ (rows/columns of all
weight matrices permuted consistently).
The model $f_\Pi$ computes an identical function to $f$, but with neurons relabelled.
We claim that the PCA subspace $\mathrm{span}(C^\ell)$ is approximately invariant under
$\Pi_\ell$:
\begin{equation}
  \text{SubspaceSim}\bigl(C^\ell,\; C^\ell_\Pi\bigr) \approx 1,
  \label{eq:perm_inv}
\end{equation}
where $\text{SubspaceSim}$ is the mean cosine of principal angles (computed via SVD of
$C^\ell (C^\ell_\Pi)^\top$).

\textbf{Intuition.}
Behavioural semantics (e.g., sycophancy, formality) are invariant to the labelling of
individual neurons.
PCA of contrastive diffs recovers the \emph{semantic} subspace, not neuron-specific
directions, hence it is $\mathcal{S}_d$-invariant.

% ======================================================================
\section{Experiments}
\label{sec:experiments}

\subsection{Setup}
\textbf{Models.}
We evaluate on four instruction-tuned models:
Llama-3.1-8B-Instruct, Qwen2.5-7B-Instruct, Gemma-2-9B-it, and Mistral-7B-Instruct-v0.3.

\textbf{Behaviors.}
Five behavioral axes with hand-crafted contrastive pairs (50--60 pairs each):
sycophancy suppression, refusal calibration, verbosity control, formality, and
uncertainty expression.

\textbf{Evaluation metric.}
\emph{Activation direction accuracy}: fraction of test pairs where the steered output's
last-token activation is cosine-closer to the positive-example activation than the
negative-example activation.
Chance level = 0.50.

\textbf{Baselines.}
(1) No steering; (2) Raw mean-diff addition (no PCA); (3) PCA directions with $K=1$ (no
calibration); (4) PCA + $K$-calibration (ours).

\subsection{Permutation Invariance}
\label{sec:exp_perm}

We apply 5 independent random neuron permutations to a random 50\% of layers and
re-extract PCA directions from the permuted model using identical prompts.
Table~\ref{tab:perm_inv} reports mean subspace cosine similarity across permutations and
layers.
All models achieve $>0.85$, substantially above the random baseline ($1/\sqrt{5} \approx
0.45$), confirming Eq.~\eqref{eq:perm_inv}.

\input{table_perm_inv}

\subsection{Weight-Space Alignment}
\label{sec:exp_weight}

For each layer, we compute the top-10 right singular vectors of $W_\mathrm{down}$ and
measure their maximum cosine alignment with each PCA component.
Figure~\ref{fig:alignment} (left) shows PCA directions align at $0.32$--$0.41$ on average,
vs.\ a random baseline of $0.10$--$0.12$, a 3--4$\times$ enrichment.
This is direct empirical evidence that behavioral PCA directions are biased toward
weight-space-canonical directions.

\begin{figure}[t]
  \centering
  \includegraphics[width=0.48\linewidth]{weight_space_alignment.pdf}
  \hfill
  \includegraphics[width=0.48\linewidth]{k_spectral_correlation.pdf}
  \caption{%
    \textbf{Left:} Heatmap of mean PCA--singular-vector alignment per layer and behavior
    (Llama-3.1-8B).
    PCA directions are consistently better aligned with weight singular vectors than random
    directions (side panel).
    \textbf{Right:} Scatter of $K$-values vs.\ spectral norms of $W_\mathrm{down}$ across
    all layers and models.
    Pearson $r > 0.70$ confirms that norm-based $K$-calibration tracks spectral geometry.
  }
  \label{fig:alignment}
\end{figure}

\subsection{$K$-Calibration Correlation}
Table~\ref{tab:k_spectral} reports Pearson and Spearman correlations between
layer-wise $K$-values and $W_\mathrm{down}$ spectral norms.
Pearson $r > 0.70$ across all models (Figure~\ref{fig:alignment}, right), supporting
Proposition~\ref{prop:k}.

\input{table_k_spectral}

\subsection{Baking Efficacy}
\label{sec:exp_efficacy}

Table~\ref{tab:efficacy} compares the four methods.
PCA + $K$-calibration outperforms raw addition by $+5$--$12$ accuracy points across
behaviors and models.
Ablating $K$ (PCA only) degrades performance by $3$--$8$ points, confirming the value of
spectral calibration.

\input{table_efficacy}

\subsection{Cross-Architecture Generalization}
Figure~\ref{fig:cka} shows CKA~\cite{kornblith2019similarity} similarity matrices between
PCA direction subspaces across architectures.
Mean CKA $> 0.60$ for sycophancy and uncertainty behaviors suggests partially universal
behavioral geometry, with variation for refusal calibration (more architecture-specific).

\begin{figure}[t]
  \centering
  \includegraphics[width=0.48\linewidth]{cross_arch_cka.pdf}
  \hfill
  \includegraphics[width=0.48\linewidth]{efficacy_comparison.pdf}
  \caption{%
    \textbf{Left:} Cross-architecture CKA similarity of PCA behavioral directions.
    High values indicate universal behavioral geometry.
    \textbf{Right:} Efficacy comparison across methods and behaviors (mean over 4 models).
    PCA + $K$-calibration (green) consistently outperforms baselines.
  }
  \label{fig:cka}
\end{figure}

% ======================================================================
\section{Conclusion}
\label{sec:conclusion}

We present \emph{activation baking}: a framework for behavioural adaptation that is
principled through weight-space symmetry analysis.
Our three contributions—PCA-directed extraction, $K$-calibration via spectral geometry,
and the permutation invariance guarantee—together establish that activation steering
operates along weight-space-canonical directions.
This reframing has practical consequences: baked behavioral directions can be composed
using the same model-merging arithmetic that governs weight-space model combination,
opening a new avenue for lightweight, composable, safety-aligned model adaptation.

\paragraph{Limitations.}
We evaluate on instruction-tuned 7--9B models; scaling to larger base models and
multi-token interventions remains future work.

\paragraph{Broader impact.}
Principled behavioral adapters lower the barrier to safe LLM customisation, which could
be used both to improve alignment and (by misuse) to instil harmful behaviours.
We release all code and contrastive-pair datasets to support responsible research.

% ======================================================================
\bibliographystyle{plainnat}
\bibliography{references}

\end{document}
